\resizebox{\textwidth}{!}{\begin{tabular}{llrrrrr}
\toprule
Surface & Measure & Cheetah fixes & Available & Expected direction & 20-km buffer & Animals\\
 & & (median) & (median) & (share) & (share) & (of 9)\\
\midrule
Final model (vegetation-balanced, documented fences) & Resistance (lower expected) & 2.23 & 2.45 & 0.60 & 0.60 & 7\\
 & Current (higher expected) & 0.156 & 0.145 & 0.65 & 0.62 & 6\\
Earlier model (pressure and terrain only, no fences) & Resistance (lower expected) & 2.03 & 2.34 & 0.61 & 0.61 & --\\
 & Current (higher expected) & 0.143 & 0.128 & 0.65 & 0.63 & --\\
\bottomrule
\end{tabular}}
\\[0.5em]
\noindent\TODO{Final note, in your own words: the comparison uses 1,006 daily locations from nine free-roaming cheetahs (2003--2008) against 10,060 random points from the same landscape, on the 2012 surfaces. ``Expected direction (share)'' is the chance that a random cheetah location had lower resistance (or higher current) than a random available point; 0.50 means no difference. ``20-km buffer'' repeats this with a wider available area. ``Animals'' counts how many of the nine had a median in the expected direction. No location fell inside a core. Add one sentence saying this is the broader comparison, and that in the stricter step check (Table~\ref{tab:limpopo_steps}) the resistance lean held but the current lean did not.}
